\documentclass{article}
\usepackage{amsmath}
\usepackage{graphicx}
\usepackage{geometry}
\geometry{a4paper, margin=1in}

\title{Subwavelength Lateral Heterogeneity: \\ Diffraction Processing and Lateral Tuning Theory}
\author{}
\date{\today}

\begin{document}

\maketitle

\section{Processing Workflow: Isolating and Analyzing Scattering}
To analyze the weak diffracted field generated by subwavelength lateral variations, the dominant coherent reflection must be systematically removed.

\begin{enumerate}
    \item \textbf{Background Removal / Coherent Subtraction:} Obtain the isolated scattered field $u_{scat}(x,t)$ by subtracting the mean coherent field $u_{avg}(x,t)$ from the full data $u_{full}(x,t)$.
    \begin{itemize}
        \item \textit{Method A (Synthetic Reference):} Subtract a forward-modeled homogeneous equivalent fracture.
        \item \textit{Method B (Data-Driven):} Apply a robust lateral low-pass filter (or spatial moving average) to estimate $u_{avg}(x,t)$, then compute $u_{scat}(x,t) = u_{full}(x,t) - u_{avg}(x,t)$. (We often use Truncated SVD for this).
    \end{itemize}
    \item \textbf{F-K Filtering:} Transform $u_{scat}$ to the $f-k$ domain. Apply a surgical mute or high-pass filter to $k_x$ to remove residual near-dc ($k_x \approx 0$) energy left over from imperfect background subtraction.
    \item \textbf{Diffraction Migration:} Perform prestack depth migration exclusively on $u_{scat}$. Rather than continuous flat events, this will image point-like diffractors at the phase boundaries of the internal blocks.
    \item \textbf{Multi-Band Spectral Analysis:} Perform the migration using both low-frequency and high-frequency bandpass filters. Low frequencies map the macro-structure; high frequencies extract the internal scattering variations.
\end{enumerate}

\section{Theoretical Framework: Fracture Grating \& Bragg Scattering}
We model the fracture as a thin layer with a mean background property and lateral perturbations. The internal structural variations act as a subwavelength spatial grating.

\subsection{Model Assumptions}
Let the horizontal spatial perturbation be periodic or quasi-periodic:
\begin{equation}
\delta m(x) = A \cos(K_0 x)
\end{equation}
where $K_0 = 2\pi/d$ is the spatial wavenumber of the blocks, $d$ is the block width, and $A$ is the contrast amplitude.

\subsection{Bragg Scattering \& Band-Limited Observation}
When an incident wave $U_{inc}(x, \omega)$ with wavenumber $k_{x,inc}$ interacts with this grating, the scattered field is generated via Bragg scattering:
\begin{equation}
k_{x,scat} = k_{x,inc} \pm K_0
\end{equation}
For the scattered energy to reach the surface receivers, it must propagate rather than decay exponentially. Therefore, the scattered wavenumber must lie within the propagating cone:
\begin{equation}
|k_{x,scat}| \le \frac{\omega}{v} = k_{medium}
\end{equation}
This imposes a fundamental bandpass filter on the observable internal spectrum. High spatial-frequency variations (very small $d$) generate large $K_0$. If $K_0 > k_{medium}$, the scattered waves are purely evanescent and undetectable in the far-field. The observed $f-k$ spectrum $U(k_x, \omega)$ is therefore the true lateral spectrum $S(k_x)$ truncated by the host medium's propagation capabilities:
\begin{equation}
|U(k_x, \omega)|^2 \approx |G(\omega, k_x)|^2 S(k_x)
\end{equation}
where $G$ accounts for acquisition geometry, wavelet spectrum, and wave propagation physics.

\section{Proposed "Lateral Tuning Law"}
Analogous to the vertical thin-bed tuning effect ($E \propto \sin(k_z \Delta z)$), we define a lateral tuning parameter:
\begin{equation}
\eta = \frac{d}{\lambda}
\end{equation}
where $\lambda$ is the dominant wavelength in the background medium.

The diffracted energy $E_{diff}(\omega)$ follows a lateral tuning relationship:
\begin{equation}
E_{diff}(\omega) \propto |\delta m(K_0)|^2 \cdot C(\eta)
\end{equation}
where $C(\eta)$ is the diffraction coupling efficiency:
\begin{itemize}
    \item \textbf{$\eta \gg 1$ (Geometric Regime):} Blocks are large enough to act as independent specular reflectors. Standard resolution applies.
    \item \textbf{$\eta \approx 0.5 - 1.0$ (Resonant Tuning Regime):} The spatial period perfectly matches the horizontal wavenumber limits of the propagating wave. This yields maximum constructive interference (Bragg resonance) back to the receivers. $C(\eta)$ approaches its peak.
    \item \textbf{$\eta < 0.5$ (Evanescent/Subwavelength Regime):} $K_0$ exceeds the radiation cone. $C(\eta)$ drops off exponentially as the scattering becomes primarily evanescent. The layer effectively homogenizes.
\end{itemize}

\section{Practical Measurable Quantities}
To invert for subwavelength structures, parameterize the data using the following metrics extracted from $u_{scat}$:
\begin{enumerate}
    \item \textbf{F-K Ridge Peaks ($K_{peak}$):} Search the $f-k$ spectrum for maximum energy ridges. If the structure is periodic, the peak $k_x$ directly estimates $K_0$, yielding $d = 2\pi/K_0$.
    \item \textbf{Frequency-Dependent Amplitude Decay:} Calculate total diffracted energy integrating over $k_x$ for various frequency bins.
    \item \textbf{Migrated Diffraction Focusing Energy:} Compute the localized energy of the migrated image $I(x,z)$ at the known fracture depth to invert for the perturbation contrast $A$.
\end{enumerate}

\end{document}
